\chapter*{Prólogo}\addcontentsline{toc}{chapter}{Prólogo}

\textit{Estructuras Aeroespaciales} nace como una edición de estudio independiente para la asignatura Teoría de Estructuras del Grado en Ingeniería Aeroespacial. Esta tercera edición conserva el objetivo de las ediciones anteriores ---aprender el Bloque I con profundidad y practicar problemas---, pero adopta una arquitectura de libro universitario: fuentes auditadas, notación explícita, bibliografía trazable, cajas de intuición y examen, referencias cruzadas e índice analítico.

El hilo conductor es
\[
\boxed{\text{cargas}\rightarrow\text{equilibrio}\rightarrow\text{esfuerzos internos}\rightarrow\text{tensiones}\rightarrow\text{deformaciones}}
\]

\begin{intuitionbox}
Las fórmulas no son el punto de partida. El punto de partida es el modelo: qué parte de la realidad conservamos, qué hipótesis aceptamos y qué ecuaciones quedan permitidas por esas hipótesis.
\end{intuitionbox}

\section*{Qué cubre esta edición}
La \BookEdition{} constituye el \textbf{Volumen I}: fundamentos, tensión y deformación, modelo monodimensional, equilibrio y leyes de esfuerzos, junto con problemas resueltos y preparación específica del Bloque I. La arquitectura queda preparada para incorporar después axil/flexión, torsión de pared delgada, cortante, métodos energéticos y elasticidad tridimensional, que aparecen en el temario oficial de la asignatura.\autocite{ule-guia-2026}

\begin{warningbox}
Este es un manual de estudio independiente. No es una publicación oficial de la Universidad de León ni de sus profesores. Los materiales de clase y exámenes de terceros se utilizan como fuentes de contraste y no se reproducen de forma sustitutiva.
\end{warningbox}
